% Options for packages loaded elsewhere
\PassOptionsToPackage{unicode}{hyperref}
\PassOptionsToPackage{hyphens}{url}
\PassOptionsToPackage{dvipsnames,svgnames,x11names}{xcolor}
\documentclass[
  a4paper,
  10pt]{article}
\usepackage{xcolor}
\usepackage[margin=22mm]{geometry}
\usepackage{amsmath,amssymb}
\setcounter{secnumdepth}{-\maxdimen} % remove section numbering
\usepackage{iftex}
\ifPDFTeX
  \usepackage[T1]{fontenc}
  \usepackage[utf8]{inputenc}
  \usepackage{textcomp} % provide euro and other symbols
\else % if luatex or xetex
  \usepackage{unicode-math} % this also loads fontspec
  \defaultfontfeatures{Scale=MatchLowercase}
  \defaultfontfeatures[\rmfamily]{Ligatures=TeX,Scale=1}
\fi
\usepackage{lmodern}
\ifPDFTeX\else
  % xetex/luatex font selection
\fi
% Use upquote if available, for straight quotes in verbatim environments
\IfFileExists{upquote.sty}{\usepackage{upquote}}{}
\IfFileExists{microtype.sty}{% use microtype if available
  \usepackage[]{microtype}
  \UseMicrotypeSet[protrusion]{basicmath} % disable protrusion for tt fonts
}{}
\makeatletter
\@ifundefined{KOMAClassName}{% if non-KOMA class
  \IfFileExists{parskip.sty}{%
    \usepackage{parskip}
  }{% else
    \setlength{\parindent}{0pt}
    \setlength{\parskip}{6pt plus 2pt minus 1pt}}
}{% if KOMA class
  \KOMAoptions{parskip=half}}
\makeatother
\usepackage{color}
\usepackage{fancyvrb}
\newcommand{\VerbBar}{|}
\newcommand{\VERB}{\Verb[commandchars=\\\{\}]}
\DefineVerbatimEnvironment{Highlighting}{Verbatim}{commandchars=\\\{\}}
% Add ',fontsize=\small' for more characters per line
\newenvironment{Shaded}{}{}
\newcommand{\AlertTok}[1]{\textcolor[rgb]{1.00,0.00,0.00}{\textbf{#1}}}
\newcommand{\AnnotationTok}[1]{\textcolor[rgb]{0.38,0.63,0.69}{\textbf{\textit{#1}}}}
\newcommand{\AttributeTok}[1]{\textcolor[rgb]{0.49,0.56,0.16}{#1}}
\newcommand{\BaseNTok}[1]{\textcolor[rgb]{0.25,0.63,0.44}{#1}}
\newcommand{\BuiltInTok}[1]{\textcolor[rgb]{0.00,0.50,0.00}{#1}}
\newcommand{\CharTok}[1]{\textcolor[rgb]{0.25,0.44,0.63}{#1}}
\newcommand{\CommentTok}[1]{\textcolor[rgb]{0.38,0.63,0.69}{\textit{#1}}}
\newcommand{\CommentVarTok}[1]{\textcolor[rgb]{0.38,0.63,0.69}{\textbf{\textit{#1}}}}
\newcommand{\ConstantTok}[1]{\textcolor[rgb]{0.53,0.00,0.00}{#1}}
\newcommand{\ControlFlowTok}[1]{\textcolor[rgb]{0.00,0.44,0.13}{\textbf{#1}}}
\newcommand{\DataTypeTok}[1]{\textcolor[rgb]{0.56,0.13,0.00}{#1}}
\newcommand{\DecValTok}[1]{\textcolor[rgb]{0.25,0.63,0.44}{#1}}
\newcommand{\DocumentationTok}[1]{\textcolor[rgb]{0.73,0.13,0.13}{\textit{#1}}}
\newcommand{\ErrorTok}[1]{\textcolor[rgb]{1.00,0.00,0.00}{\textbf{#1}}}
\newcommand{\ExtensionTok}[1]{#1}
\newcommand{\FloatTok}[1]{\textcolor[rgb]{0.25,0.63,0.44}{#1}}
\newcommand{\FunctionTok}[1]{\textcolor[rgb]{0.02,0.16,0.49}{#1}}
\newcommand{\ImportTok}[1]{\textcolor[rgb]{0.00,0.50,0.00}{\textbf{#1}}}
\newcommand{\InformationTok}[1]{\textcolor[rgb]{0.38,0.63,0.69}{\textbf{\textit{#1}}}}
\newcommand{\KeywordTok}[1]{\textcolor[rgb]{0.00,0.44,0.13}{\textbf{#1}}}
\newcommand{\NormalTok}[1]{#1}
\newcommand{\OperatorTok}[1]{\textcolor[rgb]{0.40,0.40,0.40}{#1}}
\newcommand{\OtherTok}[1]{\textcolor[rgb]{0.00,0.44,0.13}{#1}}
\newcommand{\PreprocessorTok}[1]{\textcolor[rgb]{0.74,0.48,0.00}{#1}}
\newcommand{\RegionMarkerTok}[1]{#1}
\newcommand{\SpecialCharTok}[1]{\textcolor[rgb]{0.25,0.44,0.63}{#1}}
\newcommand{\SpecialStringTok}[1]{\textcolor[rgb]{0.73,0.40,0.53}{#1}}
\newcommand{\StringTok}[1]{\textcolor[rgb]{0.25,0.44,0.63}{#1}}
\newcommand{\VariableTok}[1]{\textcolor[rgb]{0.10,0.09,0.49}{#1}}
\newcommand{\VerbatimStringTok}[1]{\textcolor[rgb]{0.25,0.44,0.63}{#1}}
\newcommand{\WarningTok}[1]{\textcolor[rgb]{0.38,0.63,0.69}{\textbf{\textit{#1}}}}
\usepackage{longtable,booktabs,array}
\newcounter{none} % for unnumbered tables
\usepackage{calc} % for calculating minipage widths
% Correct order of tables after \paragraph or \subparagraph
\usepackage{etoolbox}
\makeatletter
\patchcmd\longtable{\par}{\if@noskipsec\mbox{}\fi\par}{}{}
\makeatother
% Allow footnotes in longtable head/foot
\IfFileExists{footnotehyper.sty}{\usepackage{footnotehyper}}{\usepackage{footnote}}
\makesavenoteenv{longtable}
\setlength{\emergencystretch}{3em} % prevent overfull lines
\providecommand{\tightlist}{%
  \setlength{\itemsep}{0pt}\setlength{\parskip}{0pt}}
\usepackage{fancyhdr}
\usepackage{titlesec}
\usepackage{enumitem}
\usepackage{etoolbox}
\usepackage{xcolor}
\definecolor{reportblue}{HTML}{173A5E}
\definecolor{reportgray}{HTML}{616B75}
\pagestyle{fancy}
\fancyhf{}
\fancyhead[L]{\small\sffamily\color{reportgray}Computer Networks Laboratory}
\fancyhead[R]{\small\sffamily\color{reportgray}Lab 2: Network Layer}
\fancyfoot[C]{\small\thepage}
\setlength{\headheight}{14pt}
\renewcommand{\headrulewidth}{0.3pt}
\titleformat{\section}{\Large\bfseries\sffamily\color{reportblue}}{}{0pt}{}
\titleformat{\subsection}{\large\bfseries\sffamily\color{reportblue}}{}{0pt}{}
\titleformat{\subsubsection}{\normalsize\bfseries\sffamily}{}{0pt}{}
\titlespacing*{\section}{0pt}{13pt}{6pt}
\titlespacing*{\subsection}{0pt}{10pt}{5pt}
\titlespacing*{\subsubsection}{0pt}{8pt}{3pt}
\setlist{nosep,leftmargin=1.6em}
\setlength{\emergencystretch}{2em}
\setlength{\tabcolsep}{4pt}
\renewcommand{\arraystretch}{1.12}
\AtBeginEnvironment{longtable}{\small}
\AtBeginEnvironment{verbatim}{\small}
\makeatletter
\renewcommand{\maketitle}{%
  \begin{center}
    {\LARGE\bfseries\sffamily\color{reportblue}\@title\par}
    \vspace{5pt}
    {\small\@date\par}
  \end{center}
  \vspace{4pt}
}
\makeatother
\usepackage{bookmark}
\IfFileExists{xurl.sty}{\usepackage{xurl}}{} % add URL line breaks if available
\urlstyle{same}
\hypersetup{
  pdftitle={Lab 2: Network Layer},
  colorlinks=true,
  linkcolor={blue},
  filecolor={Maroon},
  citecolor={Blue},
  urlcolor={blue},
  pdfcreator={LaTeX via pandoc}}

\title{Lab 2: Network Layer}
\usepackage{etoolbox}
\makeatletter
\providecommand{\subtitle}[1]{% add subtitle to \maketitle
  \apptocmd{\@title}{\par {\large #1 \par}}{}{}
}
\makeatother
\subtitle{Userspace IPv4 and Distance-Vector Routing}
\author{}
\date{October 27, 2026}

\begin{document}
\maketitle

\subsection{Implementation and source
arrangement}\label{implementation-and-source-arrangement}

This lab extends the userspace Ethernet library with IPv4 transmission,
local delivery, multi-hop forwarding, and dynamic routing. The routing
metric is hop count. The implementation and CP3-CP6 experiments were
verified on \texttt{lab-comp-net}, using Linux 6.8.0, GCC 13.3.0, and
libpcap 1.10.4. The experiments were performed on October 27, 2026.

The listing uses the submission layout: evaluation programs are under
{\ttfamily\path{src/test/}}, with the required vnetUtils scripts under
{\ttfamily\path{src/test/vnetUtils/helper/}}. The original development
tree and recorded commands use {\ttfamily\path{test/}} and
{\ttfamily\path{vnetUtils/helper/}}; the submitted build and experiment
runners use the relocated paths.

{\def\LTcaptype{none} % do not increment counter
\begin{longtable}[]{@{}
  >{\raggedright\arraybackslash}p{(\linewidth - 2\tabcolsep) * \real{0.3600}}
  >{\raggedright\arraybackslash}p{(\linewidth - 2\tabcolsep) * \real{0.6400}}@{}}
\toprule\noalign{}
\begin{minipage}[b]{\linewidth}\raggedright
Source
\end{minipage} & \begin{minipage}[b]{\linewidth}\raggedright
Responsibility
\end{minipage} \\
\midrule\noalign{}
\endhead
\bottomrule\noalign{}
\endlastfoot
{\ttfamily\path{src/device.c}}, \texttt{device.h} & Register interfaces;
cache their MAC and IPv4 addresses; own separate RX and TX handles. \\
{\ttfamily\path{src/packetio.c}}, \texttt{packetio.h} & Construct and
capture Ethernet II frames; dispatch the frame callback; serialize frame
injection. \\
{\ttfamily\path{src/ip.c}}, \texttt{ip.h} & PT3: construct, validate,
deliver, and forward IPv4 datagrams; register the IP callback. \\
{\ttfamily\path{src/smtroute.c}}, \texttt{smtroute.h} & PT3: manual
routes, longest-prefix lookup, DV messages, local routes, timeout, and
garbage collection. \\
{\ttfamily\path{src/smtarp.c}}, \texttt{smtarp.h} & ARP request/reply
construction and processing; synchronized ARP cache. \\
{\ttfamily\path{src/main.c}} & Register interfaces and local routes,
dispatch frames, start receive threads, and periodically maintain and
advertise routes. \\
\end{longtable}
}

Checkpoint entrypoints are
{\ttfamily\path{src/test/cp3/send_ipv4_demo.c}} (known-MAC
transmission), {\ttfamily\path{src/test/cp4/router.c}} (interactive
routing and sending for CP4/CP5), {\ttfamily\path{src/test/cp5/run.py}}
(six-node experiment), and {\ttfamily\path{src/test/cp6/lpm.c}} with
\texttt{run.py} (LPM and two-outlet capture).
{\ttfamily\path{src/test/checkpoint_common.py}} supplies shared
topology, interaction, capture, verification, and cleanup tools.

From the repository root, build the main program and checkpoint programs
with:

\begin{Shaded}
\begin{Highlighting}[]
\NormalTok{make all cp3 cp4 cp5 cp6}
\end{Highlighting}
\end{Shaded}

The build uses
\texttt{-O2\ -Wall\ -Wextra\ -Wpedantic\ -Werror\ -pthread} and links
libpcap. Executables are placed in {\ttfamily\path{tmp/}}.
\texttt{make\ cp5} reuses {\ttfamily\path{tmp/cp4-router}}. The main
program accepts one or more interface names; the checkpoint router
additionally supports \texttt{route\ \textless{}ip\textgreater{}},
\texttt{send\ \textless{}ip\textgreater{}\ \textless{}payload\textgreater{}},
and \texttt{quit}.

\subsection{WT2 - Packet trace
analysis}\label{wt2---packet-trace-analysis}

{\def\LTcaptype{none} % do not increment counter
\begin{longtable}[]{@{}
  >{\raggedright\arraybackslash}p{(\linewidth - 2\tabcolsep) * \real{0.5500}}
  >{\raggedright\arraybackslash}p{(\linewidth - 2\tabcolsep) * \real{0.4500}}@{}}
\toprule\noalign{}
\begin{minipage}[b]{\linewidth}\raggedright
Question
\end{minipage} & \begin{minipage}[b]{\linewidth}\raggedright
Answer
\end{minipage} \\
\midrule\noalign{}
\endhead
\bottomrule\noalign{}
\endlastfoot
Which ARP Reply field equals the Sender MAC in the corresponding ARP
Request? & The \textbf{Target MAC address} (target hardware address). \\
How many IPv4 packets have DF unset? & \textbf{6}. \\
Base IPv4 and IPv6 header lengths, excluding options/extensions? &
\textbf{20 bytes} for IPv4 and \textbf{40 bytes} for IPv6. \\
\end{longtable}
}

The DF count was independently checked against
{\ttfamily\path{pcap-trace/trace.pcap}}. The matching frame numbers are
12, 17, 18, 32, 38, and 44; the corresponding Wireshark display filter
is \texttt{ip.flags.df\ ==\ 0}. In the ARP exchange at frames 33 and 34,
the Request Sender MAC and Reply Target MAC both equal
\texttt{6a:15:0a:ba:9b:7c}.

\newpage

\subsection{WT3 - Determining the next-hop MAC
address}\label{wt3---determining-the-next-hop-mac-address}

The application supplies source/destination IP addresses, protocol, and
payload to \texttt{sendIPPacket()}. It does not supply an Ethernet
destination. The IP layer calls \texttt{findRouter()} with the
destination IP, selects the longest matching valid route, and obtains
the outgoing interface and next-hop MAC from that route.

The route obtains its MAC in either of the following ways:

\begin{enumerate}
\def\labelenumi{\arabic{enumi}.}
\tightlist
\item
  \textbf{Manual routing:} \texttt{setRoutingTable()} explicitly
  receives a next-hop MAC and interface. CP3 uses this interface to
  configure a host route before sending.
\item
  \textbf{Dynamic routing:} \texttt{handleDV()} records the source MAC
  of the received Ethernet advertisement and its receiving interface.
  That neighbor is the next hop for routes learned from the
  advertisement. CP4 and CP5 therefore learn both the route and the
  link-layer address needed to use it.
\end{enumerate}

\texttt{sendIPPacket()} copies the selected MAC and calls
\texttt{sendFrame()} with EtherType \texttt{0x0800}. For a remote
destination, this MAC belongs to the immediate neighbor, while the IPv4
destination remains the final host. Each forwarding router selects its
own next-hop MAC and constructs a new Ethernet header.

The source also contains an ARP fallback for a route with an all-zero
MAC: it searches the ARP cache for the route's next-hop IP and, on a
miss, sends requests and retries. This branch requires a meaningful
next-hop IP and an active receiver to process replies. The public
manual-route API supplies a MAC, and DV learns a MAC directly; the
recorded checkpoints use these two paths and do not depend on an ARP
exchange. The ARP fallback is not described as the default path for
these experiments.

\subsection{PT3 - IPv4 sending, receiving, and
forwarding}\label{pt3---ipv4-sending-receiving-and-forwarding}

\subsubsection{Sending}\label{sending}

\texttt{sendIPPacket()} checks the payload pointer and length and the
protocol range. Payloads above 1480 bytes are rejected, so the 20-byte
header and payload fit in a 1500-byte Ethernet payload. The generated
header has version 4, IHL 5, TOS 0, identification 0, DF set, fragment
offset 0, and TTL 64. Multibyte fields use network byte order, and the
IPv4 header checksum is calculated before transmission. A failed route
lookup returns \texttt{ENETUNREACH}.

\subsubsection{Receiving and local
delivery}\label{receiving-and-local-delivery}

The frame dispatcher accepts IPv4 unicast frames addressed to the
receiving interface and ignores frames whose source is that interface's
own MAC. This prevents libpcap from feeding an outgoing frame back into
the forwarding path.

\texttt{handleIPPacket()} checks version, header length, total length,
fragmentation, and header checksum. It uses the IP total length rather
than the captured Ethernet payload length, thereby excluding Ethernet
padding from local delivery. A destination matching \textbf{any
registered interface's cached IPv4 address} is local. The registered IP
callback receives the complete IP datagram, with borrowed buffer
ownership.

\subsubsection{Forwarding}\label{forwarding}

A nonlocal unicast datagram needs a valid route and TTL greater than 1.
The handler copies the datagram, decrements TTL, recomputes the header
checksum, and sends it through the selected interface and next-hop MAC.
The source and destination IP addresses and payload are preserved.
Fragments, malformed IPv4 datagrams, TTL exhaustion, and missing routes
are dropped. Normal IPv4 drops and forwarding-send failures return 0 to
the frame receive loop so that listening continues.

\newpage

\subsection{WT4 - Distance-vector
routing}\label{wt4---distance-vector-routing}

\subsubsection{Metric, initialization, and
advertisement}\label{metric-initialization-and-advertisement}

The algorithm follows distance-vector/Bellman-Ford updates with unit
link cost. A received metric \(m\) yields candidate cost
\(\min(m+1,16)\). Metrics 0-15 are reachable; \textbf{16 means
unreachable}. Each interface contributes its own IPv4 address as a local
{\ttfamily\path{/32}} route with metric 0. This lets the experiments
measure distances to host addresses without treating a shared link
subnet as a remote host.

After registering all interfaces and initializing local routes, the
program begins reception and advertises immediately. It subsequently
calls \texttt{updateRoutingTable()} and \texttt{sendDV()} approximately
once per second. The routing table has at most 64 entries.

The wire protocol is a \textbf{custom Ethernet DV protocol}, using
EtherType \texttt{0x88B6}. It uses hop-count infinity and poison
reverse, but is not a RIPv1/RIPv2 wire-format implementation.

{\def\LTcaptype{none} % do not increment counter
\begin{longtable}[]{@{}lll@{}}
\toprule\noalign{}
DV payload offset & Size & Meaning \\
\midrule\noalign{}
\endhead
\bottomrule\noalign{}
\endlastfoot
0 & 1 byte & Version, currently 1. \\
1 & 1 byte & Reserved, must be zero. \\
2-3 & 2 bytes & Entry count, network byte order. \\
Each entry +0 & 4 bytes & Destination prefix. \\
Each entry +4 & 4 bytes & Subnet mask. \\
Each entry +8 & 4 bytes & Advertised metric, network byte order. \\
\end{longtable}
}

The header is 4 bytes and each entry is 12 bytes. A full 64-entry
advertisement occupies 772 bytes, within the Ethernet payload limit.
Advertisements use the Ethernet broadcast address on each interface.

\subsubsection{Processing an
advertisement}\label{processing-an-advertisement}

Before changing routes, \texttt{handleDV()} validates the complete
message: source MAC, version, reserved byte, count/length bounds,
contiguous masks, canonical prefixes, and metric range. Ethernet padding
is permitted after the declared entries.

Routes are identified by \textbf{destination prefix and mask}. The
receiving interface together with the Ethernet source MAC identifies the
neighbor. A new reachable prefix is inserted if capacity permits. For an
existing dynamic route:

\begin{itemize}
\tightlist
\item
  An update from the current next hop is accepted even if the metric
  becomes worse or unreachable.
\item
  A different neighbor replaces the current route when its reachable
  path is shorter, or when the current route is invalid.
\item
  Equal-cost advertisements from another neighbor retain the current
  next hop.
\item
  Local and manually configured routes are protected from dynamic
  replacement.
\end{itemize}

Only the selected path is stored. The implementation does not retain
every neighbor's alternative vector. Prefix length determines forwarding
selection; hop count compares alternatives to the \textbf{same} prefix.

\subsubsection{Poison reverse, timeout, and
cleanup}\label{poison-reverse-timeout-and-cleanup}

For each outgoing interface, \texttt{sendDV()} takes a synchronized
table snapshot. A dynamic route learned through that interface is
advertised back with metric 16. Local/manual routes are advertised
normally. Poison reverse changes the outgoing advertisement, not the
stored route.

A dynamic route expires when no accepted reachable update has refreshed
it for 5 seconds. It becomes invalid with metric 16 and remains
available for withdrawal advertisements for another 5 seconds before
deletion. Repeated unreachable updates do not restart this
garbage-collection timer. A suitable reachable update can restore an
invalid route. Lookup skips invalid/infinite routes immediately, without
waiting for deletion.

\newpage

\subsection{Concurrency and boundary
conditions}\label{concurrency-and-boundary-conditions}

\subsubsection{Reentrancy in the running
stack}\label{reentrancy-in-the-running-stack}

Each registered interface has separate RX and TX libpcap handles. One
receive thread owns each RX handle. \texttt{sendFrame()} serializes
injection and associated TX error access with a mutex, allowing periodic
DV transmission and receive-triggered forwarding to share the sending
path.

The routing table and ARP cache each have a mutex. Route updates,
timeout/cleanup, and lookup execute under the route mutex.
\texttt{findRouter()} returns a \textbf{thread-local copy}, so a caller
does not retain a pointer into a table that another thread may move
during garbage collection. The copy remains usable until that thread's
next lookup. DV sending copies a snapshot under the mutex and releases
the lock before frame transmission.

Interface registration and callback installation occur before receive
threads start. The registry is then read-only. Cleanup is performed
after its users have stopped, or by terminating the standalone
experiment process. Concurrent registry reconfiguration or callback
replacement is outside this operating contract. Application callbacks
can be invoked by different interface threads and must synchronize their
own shared application state.

\subsubsection{Conditions addressed in the
implementation}\label{conditions-addressed-in-the-implementation}

{\def\LTcaptype{none} % do not increment counter
\begin{longtable}[]{@{}
  >{\raggedright\arraybackslash}p{(\linewidth - 2\tabcolsep) * \real{0.3400}}
  >{\raggedright\arraybackslash}p{(\linewidth - 2\tabcolsep) * \real{0.6600}}@{}}
\toprule\noalign{}
\begin{minipage}[b]{\linewidth}\raggedright
Condition
\end{minipage} & \begin{minipage}[b]{\linewidth}\raggedright
Handling
\end{minipage} \\
\midrule\noalign{}
\endhead
\bottomrule\noalign{}
\endlastfoot
Ethernet padding after a short IP packet & IP total length bounds the
delivered or forwarded datagram. \\
Captured outgoing IPv4 frame & Dispatcher rejects an interface's own
source MAC and frames not addressed to its MAC. \\
Destination is another local interface & Local-address detection checks
all registered interfaces. \\
Bad IPv4 version, lengths, or checksum & Drop the datagram and continue
listening. \\
Fragmented IPv4 packet & Drop, as expressly permitted by the handout. \\
TTL at most 1 or no forwarding route & Drop without terminating the
receive loop. \\
Malformed DV entry late in a message & Validate all entries before any
route mutation. \\
Current next hop advertises a worse metric & Accept it, allowing loss of
reachability to propagate. \\
Equal-cost alternative & Keep the selected neighbor to avoid unnecessary
switching. \\
Repeated poison advertisement & Preserve the first invalidation time,
allowing eventual cleanup. \\
Routing table at capacity & Refuse new entries; existing routes can
still be updated. \\
Interface temporarily down & Log send failures; periodic processing
continues. Recovery is exercised in CP4/CP5. \\
\end{longtable}
}

The implementation uses a maximum of 16 registered devices and 64
routes, and caches the first IPv4 address of each interface at
registration. Interface addresses are configured before registration in
all experiments. The ARP cache contains up to 128 IP-to-MAC mappings and
is protected independently from the route table.

Poison reverse reduces loops, but it does not guarantee instantaneous
convergence or eliminate every multi-router count-to-infinity scenario.
Periodic updates, finite infinity, and route expiry determine eventual
behavior. The checkpoint results establish the recorded stable states
and actual deliveries, rather than a universal convergence-time bound.

All mandatory Lab 2 tasks are covered. Sending uses default TOS and no
IP options; received options are not interpreted. Dropping fragments and
omitting TOS/options support follow Section 3.3.1 of the handout.
Accordingly, \texttt{not-implemented.pdf} is an empty file, as
instructed for a completed lab.

\newpage

\subsection{CP3 - IPv4 construction and packet
bytes}\label{cp3---ipv4-construction-and-packet-bytes}

\textbf{Evidence:} {\ttfamily\path{checkpoints/cp3/capture.typescript}},
\texttt{send.typescript}, and \texttt{ipv4.pcap}.

The topology is a direct link between \texttt{cp3ns1}
({\ttfamily\path{10.33.1.1/24}}) and \texttt{cp3ns2}
({\ttfamily\path{10.33.1.2/24}}). The sender installs a
{\ttfamily\path{/32}} manual route using the peer's actual MAC and calls
\texttt{sendIPPacket()} with protocol 253 and the eight-byte string
\texttt{lab2-cp3}. This experiment verifies IPv4 construction and
transmission; dynamic multi-hop routing is checked separately below.

The complete captured Ethernet frame is:

\begin{Shaded}
\begin{Highlighting}[]
\NormalTok{0x0000: 7675 a67f 1b5f 4227 0d8e e9dd 0800 4500}
\NormalTok{0x0010: 001c 0000 4000 40fd 23a1 0a21 0101 0a21}
\NormalTok{0x0020: 0102 6c61 6232 2d63 7033 0000 0000 0000}
\NormalTok{0x0030: 0000 0000 0000 0000 0000 0000}
\end{Highlighting}
\end{Shaded}

Offsets below are \textbf{decimal and zero-based from the Ethernet frame
start}. The IPv4 datagram starts at frame offset 14; its first byte is
\texttt{0x45}.

{\def\LTcaptype{none} % do not increment counter
\begin{longtable}[]{@{}
  >{\raggedright\arraybackslash}p{(\linewidth - 4\tabcolsep) * \real{0.1200}}
  >{\raggedright\arraybackslash}p{(\linewidth - 4\tabcolsep) * \real{0.2900}}
  >{\raggedright\arraybackslash}p{(\linewidth - 4\tabcolsep) * \real{0.5900}}@{}}
\toprule\noalign{}
\begin{minipage}[b]{\linewidth}\raggedright
Frame bytes
\end{minipage} & \begin{minipage}[b]{\linewidth}\raggedright
Value
\end{minipage} & \begin{minipage}[b]{\linewidth}\raggedright
Interpretation
\end{minipage} \\
\midrule\noalign{}
\endhead
\bottomrule\noalign{}
\endlastfoot
0-5 & \texttt{76\ 75\ a6\ 7f\ 1b\ 5f} & Destination MAC:
\texttt{76:75:a6:7f:1b:5f}. \\
6-11 & \texttt{42\ 27\ 0d\ 8e\ e9\ dd} & Source MAC:
\texttt{42:27:0d:8e:e9:dd}. \\
12-13 & \texttt{08\ 00} & EtherType \texttt{0x0800}: IPv4. \\
14 & \texttt{45} & High 4 bits \texttt{0100}: version 4. Low 4 bits
\texttt{0101}: IHL 5, or 20 bytes. \\
15 & \texttt{00} & DSCP and ECN are zero. \\
16-17 & \texttt{00\ 1c} & IPv4 total length: 28 bytes. \\
18-19 & \texttt{00\ 00} & Identification: 0. \\
20-21 & \texttt{40\ 00} & Flags \texttt{010}: reserved 0, DF 1, MF 0.
The 13-bit fragment offset is 0. \\
22 & \texttt{40} & TTL: 64. \\
23 & \texttt{fd} & Protocol: 253. The test payload has no TCP/UDP
header. \\
24-25 & \texttt{23\ a1} & IPv4 header checksum: \texttt{0x23a1}. \\
26-29 & \texttt{0a\ 21\ 01\ 01} & Source address: \texttt{10.33.1.1}. \\
30-33 & \texttt{0a\ 21\ 01\ 02} & Destination address:
\texttt{10.33.1.2}. \\
34-41 & \texttt{6c\ 61\ 62\ 32\ 2d\ 63\ 70\ 33} & ASCII
\texttt{l\ a\ b\ 2\ -\ c\ p\ 3}, with no trailing NUL. \\
42-59 & Eighteen \texttt{00} bytes & Ethernet padding, outside the IPv4
datagram. \\
\end{longtable}
}

The captured frame has \(14+28+18=60\) bytes, excluding FCS. The IP
length is \(20+8=28\) bytes. Independently summing the 16-bit header
words with end-around carry, including the checksum, yields
\texttt{0xffff}, confirming a valid header checksum. Both terminal
recordings ended normally.

To inspect the original bytes from the repository root:

\begin{Shaded}
\begin{Highlighting}[]
\NormalTok{sudo tcpdump {-}nn {-}e {-}vvv {-}XX {-}r checkpoints/cp3/ipv4.pcap}
\end{Highlighting}
\end{Shaded}

\newpage

\subsection{CP4 - Four-node discovery, failure, and
recovery}\label{cp4---four-node-discovery-failure-and-recovery}

\textbf{Evidence directory:} {\ttfamily\path{checkpoints/cp4/}}.

\textbf{Recordings:} \texttt{control.typescript},
\texttt{ns1.typescript} through \texttt{ns4.typescript}.
\textbf{Capture:} \texttt{forward.pcap}.

\begin{Shaded}
\begin{Highlighting}[]
\NormalTok{ns1 {-}{-}{-}{-}{-}{-}{-}{-} ns2 {-}{-}{-}{-}{-}{-}{-}{-} ns3 {-}{-}{-}{-}{-}{-}{-}{-} ns4}
\NormalTok{       10.34.1/24   10.34.2/24   10.34.3/24}
\end{Highlighting}
\end{Shaded}

{\def\LTcaptype{none} % do not increment counter
\begin{longtable}[]{@{}lll@{}}
\toprule\noalign{}
Link & Left interface address & Right interface address \\
\midrule\noalign{}
\endhead
\bottomrule\noalign{}
\endlastfoot
ns1-ns2 & {\ttfamily\path{10.34.1.1/24}} &
{\ttfamily\path{10.34.1.2/24}} \\
ns2-ns3 & {\ttfamily\path{10.34.2.1/24}} &
{\ttfamily\path{10.34.2.2/24}} \\
ns3-ns4 & {\ttfamily\path{10.34.3.1/24}} &
{\ttfamily\path{10.34.3.2/24}} \\
\end{longtable}
}

Each namespace runs one checkpoint router with all its interfaces.
Kernel IPv4 forwarding is disabled, and \texttt{bypassKernel} installs
INPUT, OUTPUT, and FORWARD DROP rules. Packet transmission and
forwarding therefore use the userspace implementation. No manual remote
routes are added.

The test queries ns1's route to \texttt{10.34.3.2} and ns4's route to
\texttt{10.34.1.1}. The query runs \textbf{inside the same process} as
the routing protocol, since its route table is process-local.

{\def\LTcaptype{none} % do not increment counter
\begin{longtable}[]{@{}
  >{\raggedright\arraybackslash}p{(\linewidth - 6\tabcolsep) * \real{0.2200}}
  >{\raggedright\arraybackslash}p{(\linewidth - 6\tabcolsep) * \real{0.2300}}
  >{\raggedright\arraybackslash}p{(\linewidth - 6\tabcolsep) * \real{0.2300}}
  >{\raggedright\arraybackslash}p{(\linewidth - 6\tabcolsep) * \real{0.3200}}@{}}
\toprule\noalign{}
\begin{minipage}[b]{\linewidth}\raggedright
Phase
\end{minipage} & \begin{minipage}[b]{\linewidth}\raggedright
ns1 to ns4
\end{minipage} & \begin{minipage}[b]{\linewidth}\raggedright
ns4 to ns1
\end{minipage} & \begin{minipage}[b]{\linewidth}\raggedright
Data delivery
\end{minipage} \\
\midrule\noalign{}
\endhead
\bottomrule\noalign{}
\endlastfoot
Initial convergence & Metric 3, learned from DV & Metric 3, learned from
DV & \texttt{cp4-before} arrives, TTL 62. \\
Both ns2 interfaces down & \texttt{UNREACHABLE} & \texttt{UNREACHABLE} &
Send command reports no valid route. \\
Both interfaces restored & Metric 3, learned from DV & Metric 3, learned
from DV & \texttt{cp4-after} arrives, TTL 62. \\
\end{longtable}
}

The selected ns1 outlet is \texttt{cp4veth1-2}, with next-hop MAC
\texttt{d2:53:38:07:e9:55}, the adjacent ns2 interface. The actual ns4
receive records include:

\begin{Shaded}
\begin{Highlighting}[]
\NormalTok{RX src=10.34.1.1 dst=10.34.3.2 ttl=62 payload=cp4{-}before}
\NormalTok{RX src=10.34.1.1 dst=10.34.3.2 ttl=62 payload=cp4{-}after}
\end{Highlighting}
\end{Shaded}

For the failure phase, \texttt{ip\ link\ set\ ...\ down} disables both
ns2 interfaces. Its sending errors are retained in the recording. The
four processes remain running. After dynamic routes expire, both
endpoints report unreachable. Bringing the interfaces back up allows
those same processes to relearn the routes.

\texttt{forward.pcap} independently contains the two expected datagrams,
with valid IPv4 header checksums. TTL starts at 64 and becomes 62 after
forwarding at ns2 and ns3; the destination host does not decrement it
for local delivery. All five terminal recordings ended normally.

These observations satisfy the three CP4 states: discovery, loss of
discovery after isolating ns2, and recovery after reconnection. They do
not measure exact convergence latency. An unreachable lookup establishes
the absence of a usable route, rather than directly inspecting whether
its invalid entry has already been garbage-collected.

\newpage

\subsection{CP5 - Six-node distance matrix and
isolation}\label{cp5---six-node-distance-matrix-and-isolation}

\textbf{Evidence:} {\ttfamily\path{checkpoints/cp5/control.typescript}},
\texttt{cp5ns1.typescript} through \texttt{cp5ns6.typescript}, and
\texttt{ns1.pcap} through \texttt{ns6.pcap}.

\begin{Shaded}
\begin{Highlighting}[]
\NormalTok{ns1 {-}{-}{-} ns2 {-}{-}{-} ns3 {-}{-}{-} ns4}
\NormalTok{         |       |}
\NormalTok{        ns5 {-}{-}{-} ns6}
\end{Highlighting}
\end{Shaded}

The six links are 1-2, 2-3, 3-4, 2-5, 3-6, and 5-6, using
{\ttfamily\path{10.35.1.0/24}} through {\ttfamily\path{10.35.6.0/24}} in
that order. In particular, \textbf{ns6 connects to ns3}. Each router
registers all its interfaces. As in CP4, kernel IP processing and
forwarding are bypassed.

Each host is represented by one of its registered IPv4 addresses:
ns1=\texttt{10.35.1.1}, ns2=\texttt{10.35.1.2}, ns3=\texttt{10.35.2.2},
ns4=\texttt{10.35.3.2}, ns5=\texttt{10.35.4.2}, and
ns6=\texttt{10.35.5.2}. The measured entries below come from actual
\texttt{findRouter()} queries. An independent breadth-first search on
the topology supplies the expected hop counts. Three successive complete
query rounds must agree before a phase is accepted.

\subsubsection{Initial measured
distances}\label{initial-measured-distances}

{\def\LTcaptype{none} % do not increment counter
\begin{longtable}[]{@{}lrrrrrr@{}}
\toprule\noalign{}
Source / destination & ns1 & ns2 & ns3 & ns4 & ns5 & ns6 \\
\midrule\noalign{}
\endhead
\bottomrule\noalign{}
\endlastfoot
ns1 & 0 & 1 & 2 & 3 & 2 & 3 \\
ns2 & 1 & 0 & 1 & 2 & 1 & 2 \\
ns3 & 2 & 1 & 0 & 1 & 2 & 1 \\
ns4 & 3 & 2 & 1 & 0 & 3 & 2 \\
ns5 & 2 & 1 & 2 & 3 & 0 & 1 \\
ns6 & 3 & 2 & 1 & 2 & 1 & 0 \\
\end{longtable}
}

\subsubsection{Measured distances after isolating
ns5}\label{measured-distances-after-isolating-ns5}

Both ns5 interfaces, \texttt{cp5veth5-2} and \texttt{cp5veth5-6}, are
set down. No routing process is restarted. After convergence, the
measured matrix is:

{\def\LTcaptype{none} % do not increment counter
\begin{longtable}[]{@{}lrrrrrr@{}}
\toprule\noalign{}
Source / destination & ns1 & ns2 & ns3 & ns4 & ns5 & ns6 \\
\midrule\noalign{}
\endhead
\bottomrule\noalign{}
\endlastfoot
ns1 & 0 & 1 & 2 & 3 & \(\infty\) & 3 \\
ns2 & 1 & 0 & 1 & 2 & \(\infty\) & 2 \\
ns3 & 2 & 1 & 0 & 1 & \(\infty\) & 1 \\
ns4 & 3 & 2 & 1 & 0 & \(\infty\) & 2 \\
ns5 & \(\infty\) & \(\infty\) & \(\infty\) & \(\infty\) & 0 &
\(\infty\) \\
ns6 & 3 & 2 & 1 & 2 & \(\infty\) & 0 \\
\end{longtable}
}

Here \(\infty\) means no usable route, represented by protocol metric
16. The other five nodes remain connected; their shortest distances do
not change because the ns2-ns3-ns6 path remains. Restoring both ns5
interfaces, with the same processes, restores the complete initial
matrix. The three phases cover 108 measured distance entries.

\subsubsection{Actual packet delivery}\label{actual-packet-delivery}

Each phase also sends a uniquely labeled test payload for every ordered
pair of distinct hosts. There are 30 successful deliveries initially, 20
while ns5 is isolated, and 30 after restoration. The 10 ordered pairs
involving isolated ns5 report unreachable without sending. Destination
callbacks and six independent receive-side captures verify all
\textbf{80 successful deliveries}, with no duplicate test payloads.
Every captured header checksum is valid, and TTL equals \(65-d\) for a
route of \(d\) links. All capture processes report zero kernel drops.

The first run stopped capture before one buffered packet was written.
That diagnostic run remains in
{\ttfamily\path{tmp/cp5-first-run-evidence}}. The final run waits for
the capture files to contain all expected packets before stopping and
passes all checks. This is a capture-completion correction; the protocol
source was unchanged. The retained CP5 directory contains the complete
passing run.

\newpage

\subsection{CP6 - Longest-prefix
matching}\label{cp6---longest-prefix-matching}

\textbf{Evidence:} {\ttfamily\path{checkpoints/cp6/control.typescript}},
\texttt{out-a.pcap}, and \texttt{out-b.pcap}.

\begin{Shaded}
\begin{Highlighting}[]
\NormalTok{cp6ns2 {-}{-}{-}{-}{-}{-}{-}{-} cp6ns1 {-}{-}{-}{-}{-}{-}{-}{-} cp6ns3}
\NormalTok{  outlet A peer   sender         outlet B peer}
\end{Highlighting}
\end{Shaded}

The sender's outlet A is \texttt{cp6veth1-2}
({\ttfamily\path{10.36.1.1/24}}), connected to \texttt{10.36.1.2}.
Outlet B is \texttt{cp6veth1-3} ({\ttfamily\path{10.36.2.1/24}}),
connected to \texttt{10.36.2.2}. Both peer interfaces are captured
independently. Kernel IP processing is bypassed.

{\ttfamily\path{src/test/cp6/lpm.c}} first confirms that an empty route
table has no route for \texttt{203.0.113.7}. It then inserts four routes
in the deliberately mixed order \textbf{/24, /0, /32, /16}. All four
routes are present before the matching and transmission checks. Their
next-hop MACs are read from the actual peer interfaces.

{\def\LTcaptype{none} % do not increment counter
\begin{longtable}[]{@{}llll@{}}
\toprule\noalign{}
Destination & Selected prefix & Outlet & Payload \\
\midrule\noalign{}
\endhead
\bottomrule\noalign{}
\endlastfoot
\texttt{10.60.1.42} & {\ttfamily\path{10.60.1.42/32}} & A &
\texttt{cp6-host} \\
\texttt{10.60.1.43} & {\ttfamily\path{10.60.1.0/24}} & B &
\texttt{cp6-subnet} \\
\texttt{10.60.2.42} & {\ttfamily\path{10.60.0.0/16}} & A &
\texttt{cp6-network} \\
\texttt{203.0.113.7} & {\ttfamily\path{0.0.0.0/0}} & B &
\texttt{cp6-default} \\
\end{longtable}
}

For each destination, the test compares the returned prefix, mask,
interface, and next-hop MAC before calling the real
\texttt{sendIPPacket()}. The {\ttfamily\path{/32}} case must override a
matching {\ttfamily\path{/24}}, {\ttfamily\path{/16}}, and default
route. The {\ttfamily\path{/24}} case overrides {\ttfamily\path{/16}}
and default, {\ttfamily\path{/16}} overrides default, and the final case
uses the default route. Because the insertion order differs from
prefix-length order, these results establish longest matching prefix
selection rather than first/last inserted selection.

The raw query/transmission output records:

\begin{Shaded}
\begin{Highlighting}[]
\NormalTok{PASS dst=10.60.1.42 prefix=10.60.1.42 mask=255.255.255.255}
\NormalTok{PASS dst=10.60.1.43 prefix=10.60.1.0 mask=255.255.255.0}
\NormalTok{PASS dst=10.60.2.42 prefix=10.60.0.0 mask=255.255.0.0}
\NormalTok{PASS dst=203.0.113.7 prefix=0.0.0.0 mask=0.0.0.0}
\end{Highlighting}
\end{Shaded}

The excerpt omits outlet/payload fields for compactness; the original
lines are in \texttt{control.typescript}. Outlet A captures exactly
\texttt{cp6-host} and \texttt{cp6-network}; outlet B captures exactly
\texttt{cp6-subnet} and \texttt{cp6-default}. All four datagrams retain
TTL 64 because this is a sending-path outlet test. Destination/source
IPs, payloads, IPv4 header checksums, and the destination Ethernet MACs
match the selected routes. Both capture processes exit normally with
zero kernel drops.

\subsection{Evidence and reproduction}\label{evidence-and-reproduction}

Each checkpoint directory keeps only \texttt{README.md}, original
\texttt{.typescript} recordings, and \texttt{.pcap} captures. Detailed
commands and the per-checkpoint explanation are in those files. CP3/CP4
were recorded interactively; CP5/CP6 were executed by the supplied
experiment runners. The CP5/CP6 runs stopped their processes and removed
their own namespaces after verification.

Build with \texttt{make\ all\ cp3\ cp4\ cp5\ cp6}. For an automatic
rerun, use {\ttfamily\path{src/test/cp5/run.py}} or
{\ttfamily\path{src/test/cp6/run.py}} with a new result-directory
argument, under \texttt{script\ -q\ -f\ -e\ -c\ ...}, as shown in each
checkpoint README. Auxiliary CSV/JSON data and capture diagnostics go to
{\ttfamily\path{tmp/}}.

This report covers PT3, WT2-WT4, and CP3-CP6. The task definitions and
permitted simplifications follow Sections 3.3 and 3.3.1 of the supplied
assignment \texttt{README.pdf}; the report/evidence arrangement follows
Section 6. Claims about execution refer to the preserved recordings and
captured packets described above.

\end{document}
